\chapter{Reflective discussion}
\label{ch:reflection}

\begin{figure}[htbp]
  \centering
  \includegraphics[width=0.95\linewidth,height=0.65\textheight,keepaspectratio]{reflection}
  \caption[Gibb's reflective cycle for iterative MLOps planning.]{\textit{Gibb's reflective cycle} applied to iterative planning for the Lichess MLOps design.}
  \label{fig:reflection-iterations}
\end{figure}

This planning pass showed how sharply storage choices diverge from pipeline choreography once volume grows:~\cref{fig:mlops-pipeline} sequences Airflow stages, whereas~\cref{ch:storage-planning} had to freeze Parquet partitions and surrogate keys before analysts could safely join~\cref{fig:storage-star-erd}.
The star-versus-wide-table tension was formative---duplicated player columns tempted me until schema churn on hundreds of millions of rows made normalised dimensions unavoidable.
Iteratively, laptop CSV pilots collapsed under monthly Lichess drops, forcing the lake ELT story aligned with~\cref{fig:etl-flow} and immutable MinIO prefixes~\cite{lichess_open_database}.
\Cref{fig:reflection-iterations} reminds me to record each Gibbs stage when diagrams change so future retrains cite exact artefact vintages.
If I repeated the task, I would co-author Great Expectations suites beside the ingestion DAG definition, checkpoint partition conventions in git with orchestration tweaks, and run prose word-limit checks beside \texttt{latexmk}.
